\documentclass[conference]{IEEEtran}
\IEEEoverridecommandlockouts
\usepackage{cite}
\usepackage{amsmath,amssymb,amsfonts}
\usepackage{algorithmic}
\usepackage{graphicx}
\usepackage{textcomp}
\usepackage{xcolor}
\usepackage{booktabs}
\usepackage{multirow}
\usepackage{url}

\def\BibTeX{{\rm B\kern-.05em{\sc i\kern-.025em b}\kern-.08em
    T\kern-.1667em\lower.7ex\hbox{E}\kern-.125emX}}

\begin{document}

\title{Federated SCAFFOLD Learning for Pneumonia Detection Across Adult and Pediatric Chest Radiograph Populations\\}

\author{
\IEEEauthorblockN{Veerana Gouda}
\IEEEauthorblockA{\textit{cHEAL, Dept. of CSE} \\
\textit{PES University, RR Campus} \\
Bengaluru 560085, India \\
veeranagoudam551@gmail.com}
\and
\IEEEauthorblockN{Satwik K}
\IEEEauthorblockA{\textit{Dept. of CSE} \\
\textit{PES University, RR Campus} \\
Bengaluru 560085, India \\
satusatwik610@gmail.com}
\and
\IEEEauthorblockN{Satwik Hegde}
\IEEEauthorblockA{\textit{Dept. of CSE} \\
\textit{PES University, RR Campus} \\
Bengaluru 560085, India \\
satwikhegdeylp11@gmail.com}
\and
\IEEEauthorblockN{Srujan Madhu Kumar}
\IEEEauthorblockA{\textit{Dept. of CSE} \\
\textit{PES University, RR Campus} \\
Bengaluru 560085, India \\
srujanmadhukumar@gmail.com}
}

\maketitle

\begin{abstract}
Automated pneumonia detection from chest radiographs can support clinical
triage, but training a single high-quality model typically requires pooling
imaging data from multiple hospitals -- something patient privacy
regulations make difficult. Federated learning addresses this by training a
shared model across institutions without centralizing raw images, but
performance can degrade when client data distributions are highly
non-identically distributed (non-IID). This paper presents a federated
pneumonia detection system built on the SCAFFOLD algorithm with a
DenseNet121 backbone, trained across four non-IID clients: three adult
clients drawn from the RSNA Pneumonia Detection Challenge dataset with
deliberately different pneumonia prevalence (12.5\%, 24.4\%, 17.7\%), and one
pediatric client drawn from the Kermany chest X-ray dataset (73.2\%
prevalence), representing a hospital serving a markedly different patient
population. On a held-out adult test set our model achieves 78.6\% accuracy
and 80.3\% F1-score; on a held-out, patient-level-disjoint pediatric test
set it achieves 88.3\% accuracy and 91.3\% F1-score. We further show that
the strategy used to select the model's decision threshold has a large and
unequal effect across the two populations: optimizing a single pooled
validation metric can silently sacrifice one subgroup's recall for the
other's precision, and we propose selecting the threshold that maximizes
the \emph{minimum} of the two subgroups' F1-scores as a simple, effective
fix. We also report a negative result from an OpenCV-based CLAHE contrast
enhancement preprocessing step, which reduced accuracy from 80.8\% to
77.2\% on the adult-only baseline, illustrating that medical-image
preprocessing changes require empirical validation rather than assumed
benefit.
\end{abstract}

\begin{IEEEkeywords}
federated learning, SCAFFOLD, pneumonia detection, chest X-ray, DenseNet121,
non-IID data, cross-demographic generalization, medical image classification
\end{IEEEkeywords}

\section{Introduction}

Pneumonia remains a leading cause of morbidity and mortality worldwide, and
chest radiography is the primary imaging modality used for its diagnosis.
Deep convolutional neural networks have demonstrated strong performance at
classifying pneumonia from chest X-rays, but training such models
conventionally requires centralizing large volumes of labeled patient
imaging data in one location. In practice, hospitals are often unable or
unwilling to share raw patient data across institutional boundaries due to
privacy regulations and data governance policies, which limits the scale and
diversity of data available to any single training pipeline.

Federated learning (FL) offers an alternative: a shared global model is
trained by aggregating locally computed updates from multiple clients (e.g.,
hospitals), without any client's raw data ever leaving its own
infrastructure. The standard FL algorithm, Federated Averaging (FedAvg)
\cite{mcmahan2017fedavg}, performs well when client data is independently
and identically distributed (IID), but its performance degrades under
non-IID conditions, which are the norm rather than the exception in
healthcare: different hospitals see different patient populations, different
disease prevalences, and different imaging equipment. SCAFFOLD
\cite{karimireddy2020scaffold} addresses this by introducing control
variates that correct for ``client drift'' -- the tendency of each client's
local model to wander toward its own local optimum during local training
steps, which otherwise slows or destabilizes global convergence under
heterogeneous data.

In this work we build a SCAFFOLD-based federated pneumonia detection system
and use it to study a particularly challenging form of non-IID
heterogeneity: combining an adult-only patient population with a pediatric
patient population within the same federation. Chest radiographs from
children differ from those of adults in lung field proportions, bone
density, and typical pathology presentation, and the two populations in this
study additionally come from different acquisition sites and imaging
pipelines, compounding the distribution shift. Rather than simply pooling
both populations and reporting one aggregate accuracy number -- which can
mask a severe failure on the minority subgroup -- we evaluate the model on
two disjoint, held-out test sets (adult and pediatric) and report both.

Our contributions are as follows:
\begin{itemize}
    \item A SCAFFOLD-based federated pneumonia detection pipeline using a
    DenseNet121 backbone, evaluated on a deliberately non-IID 3-client adult
    federation (Section~\ref{sec:baseline}).
    \item An empirically validated negative result for OpenCV-based CLAHE
    contrast enhancement as a preprocessing step, showing that a
    plausible-sounding medical-imaging preprocessing improvement can in fact
    reduce accuracy, recall, and F1-score when not validated end-to-end
    (Section~\ref{sec:clahe}).
    \item An extension of the federation to a fourth, pediatric client
    (Kermany dataset) with a patient-level train/validation/test re-split to
    prevent leakage, and an age-stratified evaluation protocol that reports
    adult and pediatric test performance separately rather than pooled
    (Section~\ref{sec:pediatric}).
    \item An analysis showing that the choice of decision-threshold
    selection objective has a large, unequal effect on the two
    demographic subgroups, and a simple max-min threshold selection rule
    that meaningfully improves the worse-off subgroup without collapsing
    the other (Section~\ref{sec:threshold}).
\end{itemize}

\section{Related Work}

\textbf{Federated learning under non-IID data.} FedAvg \cite{mcmahan2017fedavg}
established the basic federated averaging protocol, but subsequent work
identified that client drift under non-IID data slows convergence and can
hurt final accuracy. SCAFFOLD \cite{karimireddy2020scaffold} corrects for
this using control variates -- per-client and global correction terms that
are used to adjust local gradient steps -- and provides theoretical
convergence guarantees that are independent of the degree of client
heterogeneity, unlike FedAvg.

\textbf{Deep learning for pneumonia detection.} DenseNet
\cite{huang2017densenet} is a widely used convolutional architecture for
medical image classification due to its parameter efficiency and strong
gradient flow through densely connected layers, and has been used
extensively for chest X-ray classification tasks. The RSNA Pneumonia
Detection Challenge \cite{rsna2019} released a large, radiologist-labeled
adult chest radiograph dataset that has become a standard benchmark for
pneumonia/lung-opacity detection. The Kermany dataset
\cite{kermany2018cell} provided one of the first large-scale labeled
pediatric chest X-ray datasets for pneumonia classification, sourced from a
single pediatric hospital.

\textbf{Fairness and subgroup performance in federated healthcare AI.} A
growing body of work has raised concerns that models trained on pooled or
federated healthcare data can perform unevenly across demographic
subgroups, particularly when one subgroup is a small minority of the
training data. Our work is closely related to this line of research: we
treat the pediatric population as an explicit minority client, and evaluate
it as an explicit minority test set, rather than allowing it to be absorbed
into (and potentially masked by) an aggregate metric.

\section{Datasets and Preprocessing}

\subsection{Adult population: RSNA Pneumonia Detection Challenge}

The adult portion of our data is drawn from the RSNA Pneumonia Detection
Challenge dataset \cite{rsna2019}, which provides frontal chest radiographs
in DICOM format with a binary label (pneumonia/lung opacity present or
absent). We use the full training partition of this dataset: 26{,}684
images, of which 20{,}672 (77.5\%) are NORMAL and 6{,}012 (22.5\%) are
PNEUMONIA. DICOM images are converted to grayscale PNG using min-max
normalization, with photometric-inversion handling for MONOCHROME1-encoded
studies.

\subsection{Pediatric population: Kermany chest X-ray dataset}

The pediatric portion of our data is drawn from the Kermany chest X-ray
dataset \cite{kermany2018cell}, collected from pediatric patients (ages 1--5)
at Guangzhou Women and Children's Medical Center. The dataset's own
published train/test/validation split is known to have patient-level
leakage (the same patient's images can appear in more than one split), so we
discard the original split and pool all 5{,}856 images (1{,}583 NORMAL,
4{,}273 PNEUMONIA) before performing our own patient-level split, grouping
images by the patient identifier embedded in each filename and assigning
entire patient groups to a single split. This yields 4{,}113 training
images, 595 validation images, and 1{,}148 test images, at pneumonia
prevalences of 73.2\%, 72.8\%, and 72.4\% respectively -- confirming the
patient-level split did not introduce class-balance skew across splits.

\subsection{Preprocessing pipeline}

All images, regardless of source population, pass through an identical
runtime transform: conversion to grayscale, resizing to $224\times224$,
replication to 3 channels (for compatibility with the ImageNet-pretrained
DenseNet121 backbone), conversion to tensor, and normalization using
ImageNet channel statistics ($\mu = [0.485, 0.456, 0.406]$,
$\sigma = [0.229, 0.224, 0.225]$). During training only, we additionally
apply random horizontal flipping ($p=0.3$), random rotation ($\pm7^\circ$),
random affine translation/scaling jitter, and random autocontrast
($p=0.3$) as data augmentation. Applying identical preprocessing to both
populations ensures that any performance difference we observe between
adult and pediatric test results is attributable to the underlying data,
not to inconsistent preprocessing.

\section{OpenCV-Based Contrast Enhancement: A Negative Result}
\label{sec:clahe}

Motivated by the common use of Contrast Limited Adaptive Histogram
Equalization (CLAHE) in the chest X-ray CAD literature to correct for the
low, non-uniform contrast typical of raw radiographs, we implemented a
preprocessing step combining a mild Gaussian denoise (kernel size 3,
$\sigma=0.6$) followed by CLAHE (clip limit 2.0, tile grid $8\times8$),
applied identically at training and inference time, ahead of the existing
grayscale/resize/normalize pipeline. We trained an otherwise identical model
(3-client adult-only federation, 20 rounds, seed 42) with this preprocessing
step added, to isolate preprocessing as the only variable.

\begin{table}[htbp]
\caption{Effect of CLAHE + denoise preprocessing (adult-only, 3-client
federation, identical hyperparameters otherwise)}
\begin{center}
\begin{tabular}{lcc}
\toprule
\textbf{Metric} & \textbf{Baseline (no CLAHE)} & \textbf{With CLAHE + denoise} \\
\midrule
Accuracy  & 80.79\% & 77.17\% \\
Precision & 80.77\% & 86.30\% \\
Recall    & 80.83\% & 64.58\% \\
F1-score  & 80.80\% & 73.88\% \\
\bottomrule
\end{tabular}
\label{tab:clahe}
\end{center}
\end{table}

As shown in Table~\ref{tab:clahe}, CLAHE preprocessing increased precision
but sharply reduced recall (missed pneumonia cases nearly doubled, from 230
to 425 out of 1{,}200), and reduced overall accuracy and F1-score. We
attribute this to a distribution shift: CLAHE's local contrast stretching
alters pixel intensity statistics in a way that falls outside what the
ImageNet-pretrained backbone's early convolutional layers were trained to
expect, and with only 20 rounds of federated training and selective
layer unfreezing, the network does not receive enough gradient updates to
fully re-adapt to the new input distribution. We report this as a negative
result and revert to the unmodified preprocessing pipeline for all
subsequent experiments in this paper, underscoring that plausible
domain-motivated preprocessing changes in medical imaging pipelines require
empirical end-to-end validation rather than being assumed beneficial.

\section{Adult-Only Federated Baseline}
\label{sec:baseline}

\subsection{Federation design}

Our baseline federation consists of three clients, each holding a disjoint
partition of the RSNA adult dataset, deliberately constructed with different
pneumonia prevalences to create realistic non-IID conditions
(Table~\ref{tab:adult-clients}).

\begin{table}[htbp]
\caption{Adult client data distribution (training partition)}
\begin{center}
\begin{tabular}{lcccc}
\toprule
\textbf{Client} & \textbf{NORMAL} & \textbf{PNEUMONIA} & \textbf{Total} & \textbf{Prevalence} \\
\midrule
client\_0 & 7{,}469 & 1{,}070 & 8{,}539 & 12.5\% \\
client\_1 & 4{,}979 & 1{,}605 & 6{,}584 & 24.4\% \\
client\_2 & 6{,}224 & 1{,}337 & 7{,}561 & 17.7\% \\
\midrule
Total     & 18{,}672 & 4{,}012 & 22{,}684 & 17.7\% \\
\bottomrule
\end{tabular}
\label{tab:adult-clients}
\end{center}
\end{table}

An additional 1{,}600 images (800 NORMAL, 800 PNEUMONIA) are held out as a
validation set used for round selection and decision-threshold tuning, and
2{,}400 images (1{,}200 NORMAL, 1{,}200 PNEUMONIA) are held out as a final
test set, never seen during training or model selection. Together with the
training partition, this accounts for the full 26{,}684-image RSNA training
set.

\subsection{Training configuration}

The model backbone is DenseNet121, initialized with ImageNet-pretrained
weights, with only later layers unfrozen for fine-tuning. Federated training
runs for 20 rounds with 2 local epochs per client per round, batch size 16,
and a step-decayed learning rate ($3\times10^{-4} \rightarrow
2\times10^{-4} \rightarrow 1\times10^{-4}$ over the course of training). The
final model is selected from the round with the best validation F1-score
(rather than simply the last round, since later rounds are not guaranteed to
improve on a non-IID federation), and the decision threshold is chosen by
sweeping 91 candidate thresholds against the validation set.

\subsection{Results}

The baseline model achieves 80.79\% accuracy, 80.77\% precision, 80.83\%
recall, and 80.80\% F1-score on the held-out adult test set, with confusion
matrix $\begin{bmatrix}969 & 231\\230 & 970\end{bmatrix}$ (rows: actual
NORMAL/PNEUMONIA; columns: predicted NORMAL/PNEUMONIA). This is the
reference baseline against which all subsequent adult-side results in this
paper are compared.

\section{Extending the Federation to a Pediatric Client}
\label{sec:pediatric}

\subsection{Motivation}

A model trained exclusively on adult radiographs cannot be assumed to
generalize to pediatric patients: pediatric chest radiographs differ in lung
field geometry, thoracic proportions, and typical disease presentation, and
in our case additionally originate from an entirely different acquisition
site and imaging pipeline than the RSNA data. To test and improve
cross-demographic generalization within the federated setting, we add a
fourth client built from the pediatric Kermany training partition described
in Section~III, alongside the three existing adult clients
(Table~\ref{tab:all-clients}).

\begin{table}[htbp]
\caption{All four federated clients (training partition)}
\begin{center}
\begin{tabular}{lccccc}
\toprule
\textbf{Client} & \textbf{Source} & \textbf{NORMAL} & \textbf{PNEUMONIA} & \textbf{Total} & \textbf{Prev.} \\
\midrule
client\_0 & Adult (RSNA)     & 7{,}469 & 1{,}070 & 8{,}539 & 12.5\% \\
client\_1 & Adult (RSNA)     & 4{,}979 & 1{,}605 & 6{,}584 & 24.4\% \\
client\_2 & Adult (RSNA)     & 6{,}224 & 1{,}337 & 7{,}561 & 17.7\% \\
client\_3 & Pediatric (Kermany) & 1{,}104 & 3{,}009 & 4{,}113 & 73.2\% \\
\bottomrule
\end{tabular}
\label{tab:all-clients}
\end{center}
\end{table}

client\_3 is deliberately the most extreme client in the federation, both in
pneumonia prevalence (73.2\% versus 12.5--24.4\% for the adult clients) and
in the underlying patient population -- a realistic stand-in for a new
hospital with a substantially different patient mix joining an existing
federation. The validation set is correspondingly extended by pooling the
adult validation set (1{,}600 images) with a pediatric validation set
(595 images, from the same patient-level split described in Section~III),
for a combined 2{,}195-image validation set used for round and threshold
selection. Critically, the two test sets -- adult (2{,}400 images) and
pediatric (1{,}148 images) -- are kept \emph{separate} rather than pooled,
so that subgroup performance is directly visible rather than averaged away.

\subsection{Training and initial results}

Using the same SCAFFOLD configuration, hyperparameters, and DenseNet121
backbone as the adult-only baseline (20 rounds, 2 local epochs, batch size
16, seed 42), the best round was selected at round 12 by validation F1. With
the original threshold-selection objective (an F-beta metric with
$\beta=0.5$ that weights precision over recall, chosen for the adult-only
baseline to minimize false alarms), the resulting threshold of 0.36 gives
the results in Table~\ref{tab:initial-allages}.

\begin{table}[htbp]
\caption{All-ages model, precision-favoring threshold (0.36)}
\begin{center}
\begin{tabular}{lcc}
\toprule
\textbf{Metric} & \textbf{Adult test} & \textbf{Pediatric test} \\
\midrule
Accuracy  & 79.08\% & 80.66\% \\
Precision & 77.79\% & \textbf{99.84\%} \\
Recall    & 81.42\% & \textbf{73.41\%} \\
F1-score  & 79.56\% & 84.60\% \\
\bottomrule
\end{tabular}
\label{tab:initial-allages}
\end{center}
\end{table}

Adding the pediatric client costs a small amount of adult-side performance
relative to the adult-only baseline (accuracy $-1.7$pp, F1 $-1.2$pp), an
expected consequence of splitting model capacity across a more diverse
federation. More importantly, the pediatric subgroup shows an extreme
precision/recall imbalance: 99.84\% precision (only 1 false alarm in 317
healthy pediatric images) but only 73.41\% recall (221 of 831 pediatric
pneumonia cases missed). For a screening tool, a missed-disease error
(false negative) is generally more clinically costly than a false alarm, so
this imbalance represents a real limitation that motivated the analysis in
the next section.

\section{Threshold Selection Across Demographic Subgroups}
\label{sec:threshold}

\subsection{The pooled-metric threshold masks subgroup imbalance}

The threshold of 0.36 above was chosen by sweeping candidate thresholds
against the \emph{pooled} adult+pediatric validation set using an
F-beta ($\beta=0.5$) objective that explicitly favors precision over
recall. To test whether the recall problem was specifically caused by this
precision-favoring objective, we re-swept the threshold on the same trained
model and the same pooled validation set, this time using plain F1 (equal
weight on precision and recall) as the objective, without any retraining.

\begin{table}[htbp]
\caption{All-ages model, pooled-F1 threshold (0.08)}
\begin{center}
\begin{tabular}{lcc}
\toprule
\textbf{Metric} & \textbf{Adult test} & \textbf{Pediatric test} \\
\midrule
Accuracy  & 77.88\% & \textbf{93.21\%} \\
Precision & \textbf{71.76\%} & 97.36\% \\
Recall    & 91.92\% & 93.14\% \\
F1-score  & 80.60\% & \textbf{95.20\%} \\
\bottomrule
\end{tabular}
\label{tab:f1-allages}
\end{center}
\end{table}

This threshold (0.08) essentially solved the pediatric recall problem
(73.4\% $\rightarrow$ 93.1\%) and even improved pediatric F1 to 95.2\%.
However, it did so by swinging to an extreme: adult precision collapsed
from 77.79\% to 71.76\%, since a single global pooled metric let the
subgroup whose F1-vs-threshold curve was steeper (pediatric) dominate the
threshold choice, at the direct expense of the other subgroup. This
illustrates a general risk in federated systems serving heterogeneous
populations: optimizing one pooled metric on a combined validation set does
not guarantee balanced treatment of the populations that make up that pool.

\subsection{A balanced, max-min threshold selection rule}

To avoid letting either subgroup dominate, we instead compute adult-only and
pediatric-only validation F1-score independently at each candidate
threshold (using the original, unpooled adult and pediatric validation
folders), and select the threshold that maximizes the \emph{minimum} of the
two -- i.e., the threshold with the best worst-case outcome across both
populations, again without any retraining of the underlying model.

\begin{table}[htbp]
\caption{All-ages model, balanced max-min threshold (0.18)}
\begin{center}
\begin{tabular}{lcc}
\toprule
\textbf{Metric} & \textbf{Adult test} & \textbf{Pediatric test} \\
\midrule
Accuracy  & 78.63\% & 88.33\% \\
Precision & 74.52\% & 99.02\% \\
Recall    & 87.00\% & 84.72\% \\
F1-score  & 80.28\% & 91.31\% \\
\bottomrule
\end{tabular}
\label{tab:balanced-allages}
\end{center}
\end{table}

Table~\ref{tab:balanced-allages-summary} summarizes all three
threshold-selection strategies side by side. The balanced (max-min)
threshold recovers most of the adult precision lost under the pooled-F1
threshold (71.76\% $\rightarrow$ 74.52\%) while keeping pediatric recall at
a much safer 84.72\% (compared to the original 73.41\%) -- a middle ground
that does not sacrifice either population as severely as the two single-metric
strategies.

\begin{table*}[htbp]
\caption{Summary comparison of threshold-selection strategies}
\begin{center}
\begin{tabular}{lccccccccc}
\toprule
\multirow{2}{*}{\textbf{Threshold}} & \multirow{2}{*}{\textbf{Selection objective}} &
\multicolumn{4}{c}{\textbf{Adult test}} & \multicolumn{4}{c}{\textbf{Pediatric test}} \\
\cmidrule(lr){3-6} \cmidrule(lr){7-10}
 & & Acc. & Prec. & Rec. & F1 & Acc. & Prec. & Rec. & F1 \\
\midrule
0.36 & Pooled F-beta ($\beta=0.5$, precision-favoring) & 79.08\% & 77.79\% & 81.42\% & 79.56\% & 80.66\% & 99.84\% & 73.41\% & 84.60\% \\
0.08 & Pooled F1 (balanced weight, pooled sample) & 77.88\% & 71.76\% & 91.92\% & 80.60\% & 93.21\% & 97.36\% & 93.14\% & 95.20\% \\
0.18 & Max-min of adult/pediatric F1 (proposed) & 78.63\% & 74.52\% & 87.00\% & 80.28\% & 88.33\% & 99.02\% & 84.72\% & 91.31\% \\
\bottomrule
\end{tabular}
\label{tab:balanced-allages-summary}
\end{center}
\end{table*}

The confusion matrices for the balanced threshold configuration are
$\begin{bmatrix}843 & 357\\156 & 1044\end{bmatrix}$ for the adult test set
and $\begin{bmatrix}310 & 7\\127 & 704\end{bmatrix}$ for the pediatric test
set (rows: actual NORMAL/PNEUMONIA; columns: predicted NORMAL/PNEUMONIA).

\section{Discussion}

\subsection{Threshold strategy matters as much as model architecture}

The three configurations in Table~\ref{tab:balanced-allages-summary} use
\emph{the exact same trained model weights} -- only the decision threshold
differs. Yet pediatric recall ranges from 73.41\% to 93.14\% depending
purely on how that one threshold is chosen. This suggests that, in
federated systems serving multiple demographic subgroups, reporting a
subgroup-aware threshold-selection strategy is at least as important a
design decision as the choice of aggregation algorithm, and should be
validated and reported explicitly rather than left as an implicit default.

\subsection{Limitations}

The pediatric client in this study is drawn from a single institution
(Guangzhou Women and Children's Medical Center), so age effects are
confounded with institution/acquisition-device effects; we cannot fully
attribute the adult/pediatric performance gap to age alone. The pediatric
client is also substantially smaller (4{,}113 training images) than any
individual adult client, which may limit how much the global model can
specialize to pediatric-specific features under FedAvg-style
sample-weighted aggregation. Finally, none of the configurations in this
paper reach the originally targeted 90\% accuracy / 85\% precision/recall/F1
goals simultaneously for the adult subgroup, indicating that federating
across a demographically extreme client currently comes at a real,
non-trivial cost to the majority population's specialization, rather than
being cost-free.

\subsection{Future work}

Promising directions include: incorporating additional pediatric data
sources from different institutions to disentangle age effects from
site effects; per-client-aware aggregation weighting schemes that can trade
off subgroup performance more explicitly than uniform sample-count
weighting; and formal fairness-constrained threshold or loss functions that
directly optimize for balanced subgroup recall during training rather than
via post-hoc threshold selection.

\section{Conclusion}

We presented a SCAFFOLD-based federated pneumonia detection system and used
it to study cross-demographic generalization by extending a 3-client adult
federation with a fourth pediatric client. Evaluating the two populations
separately, rather than pooling them into one aggregate metric, revealed a
severe pediatric recall deficit under the federation's original
precision-favoring threshold, which we showed can be substantially mitigated
(73.4\% $\rightarrow$ 84.7\% recall) without any retraining, simply by
selecting a decision threshold that maximizes the worse-off subgroup's
F1-score rather than a single pooled metric. We also reported a negative
result for CLAHE-based contrast enhancement preprocessing, which reduced
rather than improved accuracy in our pipeline. Together, these results argue
for treating both preprocessing changes and threshold-selection strategies
as empirical questions requiring subgroup-aware validation in federated
medical imaging systems, rather than assumptions carried over from
single-population practice.

\begin{thebibliography}{00}

\bibitem{mcmahan2017fedavg}
H.~B. McMahan, E.~Moore, D.~Ramage, S.~Hampson, and B.~A.~y Arcas,
``Communication-efficient learning of deep networks from decentralized
data,'' in \emph{Proc. 20th Int. Conf. Artificial Intelligence and
Statistics (AISTATS)}, 2017, pp. 1273--1282.

\bibitem{karimireddy2020scaffold}
S.~P. Karimireddy, S.~Kale, M.~Mohri, S.~J. Reddi, S.~U. Stich, and
A.~T. Suresh, ``SCAFFOLD: Stochastic controlled averaging for federated
learning,'' in \emph{Proc. 37th Int. Conf. Machine Learning (ICML)}, 2020,
pp. 5132--5143.

\bibitem{huang2017densenet}
G.~Huang, Z.~Liu, L.~van~der Maaten, and K.~Q. Weinberger, ``Densely
connected convolutional networks,'' in \emph{Proc. IEEE Conf. Computer
Vision and Pattern Recognition (CVPR)}, 2017, pp. 4700--4708.

\bibitem{rsna2019}
Radiological Society of North America, ``RSNA Pneumonia Detection
Challenge,'' Kaggle, 2019. [Online]. Available:
\url{https://www.kaggle.com/c/rsna-pneumonia-detection-challenge}

\bibitem{kermany2018cell}
D.~S. Kermany \emph{et al.}, ``Identifying medical diagnoses and treatable
diseases by image-based deep learning,'' \emph{Cell}, vol. 172, no. 5,
pp. 1122--1131, 2018.

\end{thebibliography}

\end{document}
